\documentclass[11pt,a4paper]{article}

\usepackage[T1]{fontenc}
\usepackage[utf8]{inputenc}
\usepackage[french,english]{babel}
\usepackage{geometry}
\usepackage{xcolor}
\usepackage{tcolorbox}
\usepackage{enumitem}
\usepackage{setspace}
\usepackage{titlesec}
\usepackage{fancyhdr}
\usepackage{microtype}
\usepackage{parskip}

\geometry{
  margin=1.6cm,
  top=2.1cm,
  bottom=2.1cm
}

\definecolor{paper}{HTML}{F6EFE4}
\definecolor{ink}{HTML}{2A241F}
\definecolor{ruleline}{HTML}{D8C9B4}
\definecolor{lavender}{HTML}{C9B8E8}
\definecolor{plum}{HTML}{6B4A78}
\definecolor{rose}{HTML}{E8B4C8}
\definecolor{mint}{HTML}{C8DCC0}
\definecolor{noteblue}{HTML}{4A6FA5}
\definecolor{softyellow}{HTML}{F7E7A8}

\pagecolor{paper}
\color{ink}
\pagestyle{fancy}
\fancyhf{}
\renewcommand{\headrulewidth}{0pt}
\fancyfoot[C]{\color{plum}\small journal de fran\c{c}ais \quad\textbullet\quad \thepage}

\titleformat{\section}
  {\Large\bfseries\color{plum}}
  {}{0pt}{}[\vspace{2pt}{\color{lavender}\titlerule[2pt]}]

\titleformat{\subsection}
  {\large\bfseries\color{noteblue}}
  {}{0pt}{}

\tcbset{
  boxrule=0pt,
  arc=8pt,
  left=10pt,
  right=10pt,
  top=8pt,
  bottom=8pt,
  fonttitle=\bfseries
}

\newtcolorbox{rulebox}{
  colback=lavender!35,
  colframe=plum,
  title=formule
}

\newtcolorbox{exbox}{
  colback=white!70!paper,
  colframe=rose,
  boxrule=1pt
}

\newtcolorbox{enbox}{
  colback=mint!40,
  colframe=mint,
  boxrule=0pt
}

\newcommand{\pair}[2]{%
  \item \textbf{#1}\\
  {\color{noteblue}\textit{#2}}%
}

\newcommand{\vocab}[2]{%
  \textbf{#1} {\color{noteblue}\textit{(#2)}}%
}

\begin{document}
\selectlanguage{french}

\begin{center}
  {\color{plum}\LARGE\bfseries C'est, ce sont}\\[4pt]
  {\color{lavender}\rule{0.55\textwidth}{2.5pt}}\\[6pt]
  {\small coloriage de notes \textbullet{} goodnotes $\rightarrow$ \LaTeX}
\end{center}

\vspace{0.6em}

\begin{exbox}
\begin{itemize}[leftmargin=1.2em,itemsep=0.55em]
  \pair{Est-ce ta trousse ? --- Oui, c'est la mienne.}
       {Is this your pencil case? --- Yes, it's mine.}
  \pair{Ce sont mes livres.}
       {These are my books. \quad \vocab{livres}{books}}
  \pair{C'est ma maison.}
       {This is my house.}
  \pair{Les ciseaux, c'est pratique.}
       {\vocab{ciseaux}{scissors}, \vocab{pratique}{practical}}
  \pair{Et leurs anniversaires ? --- C'est bient\^{o}t.}
       {\vocab{anniversaires}{birthdays} \quad \vocab{bient\^{o}t}{soon / adverb}}
\end{itemize}
\end{exbox}

\begin{rulebox}
\begin{itemize}[leftmargin=1.2em,itemsep=0.4em]
  \item \texttt{c'est} + adverbe
  \item \texttt{ce sont} + adverbe + nom \quad {\small\color{ink}(le nom p\`{e}se plus que l'adverbe)}
\end{itemize}
\end{rulebox}

\begin{exbox}
\begin{itemize}[leftmargin=1.2em,itemsep=0.55em]
  \pair{Ce sont bient\^{o}t les vacances !}
       {The holidays are coming soon!}
  \pair{Ce sont les amis de mon fr\`{e}re.}
       {These are my brother's friends.}
  \pair{\`{A} qui sont ces livres ?}
       {Whose books are these?}
\end{itemize}
\end{exbox}

\newpage

\section*{celui, celle, ceux, celles}

\begin{center}
{\large
\begin{tabular}{cccc}
\textbf{celui} & \textbf{celle} & \textbf{ceux} & \textbf{celles}\\
{\small masc.\ sg.} & {\small f\'{e}m.\ sg.} & {\small masc.\ pl.} & {\small f\'{e}m.\ pl.}
\end{tabular}}
\end{center}

\vspace{0.4em}

\begin{exbox}
\begin{itemize}[leftmargin=1.2em,itemsep=0.55em]
  \pair{\`{A} qui sont ces balles ? --- Ce sont les miennes.}
       {Whose balls are these? --- They are mine.}
  \pair{Eva aime mes biscuits mais d\'{e}teste ceux de sa tante.}
       {Eva likes my biscuits but hates her aunt's.}
\end{itemize}
\end{exbox}

\section*{penser quoi}

\begin{exbox}
\begin{itemize}[leftmargin=1.2em,itemsep=0.55em]
  \pair{Oh ! j'ai regard\'{e} \c{c}a aussi hier.}
       {Oh! I watched that too yesterday.}
  \pair{Les lasagnes ? J'adore \c{c}a !}
       {Lasagna? I love it!}
  \pair{Les probabilit\'{e}s, \c{c}a m'\'{e}nerve !}
       {Probabilities get on my nerves!}
\end{itemize}
\end{exbox}

\section*{c'est beau}

\begin{exbox}
\begin{itemize}[leftmargin=1.2em,itemsep=0.55em]
  \pair{C'est beau !}
       {It's beautiful!}
  \pair{C'est vraiment beau !}
       {It's really beautiful!}
\end{itemize}
\end{exbox}

\begin{enbox}
\textbf{d\'{e}range} --- \textit{bothers me}\\[0.3em]
on fume \`{a} c\^{o}t\'{e} de moi \quad/\quad quand on fume
\end{enbox}

\vspace{0.5em}

\begin{exbox}
\begin{itemize}[leftmargin=1.2em,itemsep=0.45em]
  \item \vocab{tout le monde}{everyone}
  \item tout le monde a une t\'{e}l\'{e}
  \item tout le monde est gentil
\end{itemize}
\end{exbox}

\newpage

\section*{il y a + nom}

\begin{rulebox}
\texttt{il y a} + \texttt{du / de la / de l' / des} + nom
\end{rulebox}

\begin{exbox}
\begin{itemize}[leftmargin=1.2em,itemsep=0.55em]
  \pair{Il y a du soleil.}{sunny \quad {\small\color{plum}soleil $\sim$ pomme}}
  \pair{Il y a du vent.}{windy \quad {\small\color{plum}close the vent}}
  \pair{Il y a du brouillard.}{foggy}
  \pair{Il y a de l'orage.}{stormy}
  \pair{Il y a des nuages.}{cloudy}
  \pair{Il y a de la pluie aujourd'hui.}{rainy}
\end{itemize}
\end{exbox}

\begin{enbox}
le temps est ensoleill\'{e} \quad --- \quad \textit{it is sunny}\\[0.35em]
aussi : il fait du soleil \quad/\quad il fait soleil
\end{enbox}

\section*{\^{e}tre d'accord}

\begin{center}
{\large \textbf{\^{e}tre} + d'accord \quad {\color{noteblue}\textit{to agree}}}
\end{center}

\begin{exbox}
\begin{itemize}[leftmargin=1.2em,itemsep=0.55em]
  \pair{Je suis compl\`{e}tement d'accord avec toi.}
       {I completely agree with you.}
  \pair{Elle est toujours d'accord avec son fr\`{e}re.}
       {She always agrees with her brother.}
  \pair{Vous \^{e}tes d'accord ?}
       {Do you agree?}
  \pair{Nous sommes d'accord avec votre proposition.}
       {We agree with your offer. \quad \vocab{sommes}{\^{e}tre}}
\end{itemize}
\end{exbox}

\newpage

\section*{je voudrais / je veux}

\begin{exbox}
\begin{itemize}[leftmargin=1.2em,itemsep=0.55em]
  \pair{Je veux une glace !}
       {I want an ice cream.}
  \pair{Je voudrais un coca, s'il te pla\^{\i}t.}
       {I would like a coke, please.}
  \pair{Je veux un c\^{a}lin.}
       {I want a hug. \quad {\small\color{plum}direct}}
  \pair{Je voudrais un c\^{a}lin.}
       {I would like a hug. \quad {\small\color{plum}poli}}
  \pair{J'aimerais\ldots}
       {I would like\ldots}
\end{itemize}
\end{exbox}

\begin{rulebox}
le pr\'{e}sent \quad vs \quad le conditionnel\\[0.3em]
{\small\texttt{je veux} = I want \qquad \texttt{je voudrais} = I would like}
\end{rulebox}

\vspace{0.7em}

\section*{s'il te pla\^{\i}t}

\begin{enbox}
\textbf{s'il te pla\^{\i}t} = please\\[0.35em]
verbe \textbf{vouloir} $\rightarrow$ \texttt{veux} ou \texttt{voudrais}\\
\textit{I would want / like}
\end{enbox}

\section*{avoir peur de}

\begin{center}
{\large \textbf{avoir peur de} \quad {\color{noteblue}\textit{to be afraid / scared of}}}
\end{center}

\begin{exbox}
\begin{itemize}[leftmargin=1.2em,itemsep=0.55em]
  \pair{Elle a peur des araign\'{e}es.}
       {She is afraid of spiders.}
  \pair{J'ai peur de tomber.}
       {I am afraid of falling.}
  \pair{Je ne sais pas pourquoi.}
       {I don't know why.}
  \pair{J'ai peur de lui.}
       {I am afraid of him.}
  \pair{Elle a peur de M.\ Potier.}
       {She is afraid of Mr.\ Potier.}
  \pair{Il a peur d'Arthur.}
       {He is afraid of Arthur.}
\end{itemize}
\end{exbox}

\newpage

\section*{avoir peur de --- exemples}

\begin{exbox}
\begin{itemize}[leftmargin=1.2em,itemsep=0.55em]
  \pair{Tu as peur du noir ?}
       {Are you afraid of the dark?}
  \pair{Ils ont peur de la m\'{e}chante sorci\`{e}re.}
       {They are afraid of the wicked witch.}
  \pair{Mon chien a peur de l'orage.}
       {My dog is scared of the thunder.}
  \pair{Mon fr\`{e}re a peur des chiens.}
       {My brother is scared of dogs.}
  \pair{J'ai peur de me perdre.}
       {I'm scared of getting lost.\\
        \vocab{perdre}{lost / losing your way}}
\end{itemize}
\end{exbox}

\vspace{0.4em}

\begin{exbox}
\begin{itemize}[leftmargin=1.2em,itemsep=0.55em]
  \pair{Vous avez peur d'\'{e}chouer.}
       {You are afraid to fail.}
  \pair{Ils ont peur de prendre le train tous seuls.}
       {They are scared of taking the train on their own.}
\end{itemize}
\end{exbox}

\begin{enbox}
\vocab{prendre le train}{taking / boarding the train}\\[0.25em]
\vocab{tous seuls}{on their own}\\[0.25em]
\vocab{toujours}{always}\\[0.25em]
\vocab{nos voisins}{our neighbours}\\[0.25em]
nous avons souvent peur d'eux \quad \vocab{souvent}{often}
\end{enbox}

\vspace{1.2em}

\begin{center}
{\color{plum}\small fin du cahier \textbullet{} notes d'ao\^{u}t 2026}
\end{center}

\end{document}
